\section{硅谷体感：模型进步窗口重新洗牌}

广密的体感是，过去一个季度模型水平提升的幅度超过 2025 年全年，Opus 4.5 到 4.6 这一类跃迁让模型从 chat 问答进入真正 agentic 模式。访谈把硅谷“御三家”放在一个动态窗口里看：OpenAI、Anthropic、Gemini 各自有窗口期，今天的胜利秘籍可能成为下一阶段的毒药。

\begin{figure}[H]
\centering
\includegraphics[width=0.96\textwidth]{frontier-lab-map.png}
\caption{访谈观点中的硅谷模型公司位置图。}
\end{figure}

\begin{knowledgebox}{读图：这不是客观排名，而是路径依赖图}
图中整理的是访谈观点：Anthropic 在 Coding/Agent 上窗口好；OpenAI 的 ToC 成功可能拖慢 Coding 战略；Gemini 资源和生态强但声势不足；Meta 被视为挑战者；xAI 受战略摇摆和团队问题影响。读者应看路径依赖，而不是把它当固定排名。
\end{knowledgebox}

\subsection{Anthropic：技术细节和文化基因}

访谈中，Anthropic 被认为抓住了 Coding/Agent 的窗口。它不是 day 1 就想清楚所有事，但创始人 hands-on 看数据、重视技术细节、文化面试严格，这些使它更容易在新范式出现时组织资源。这里的重点是：文化不是口号，而会影响团队是否重视某类数据、某类产品和某类能力。

\subsection{OpenAI：过去胜利可能成为毒药}

OpenAI 在 ChatGPT 的 ToC 成功，反而可能让组织长期专注聊天入口和消费者产品，从而忽视 Coding。访谈中的判断是，过去时代的胜利秘诀可能成为下个时代的毒药。对模型公司来说，战略惯性和组织奖励机制可能比单点技术短板更危险。

\subsection{Gemini、Meta 与 xAI}

Gemini 被评价为模型能力强、生态位好，但 Coding 严重落后，Google 的战略失误在于没有更早把 Coding 作为主线。Meta 被视为新挑战者，取代 xAI 成为四号种子；xAI 则因战略摇摆和 founding team 流失被认为短期掉队。这里反映的是：前沿模型竞争已经不只是模型参数，而是组织、战略、产品和执行力。

\begin{importantbox}{硅谷御三家的真实问题}
领先模型公司不重视 Coding，就可能掉出第一梯队。因为 Coding 是模型进入数字世界的执行接口，也是加速自身研究的核心工具。
\end{importantbox}

\subsection{本章小结}

模型公司窗口在快速轮换。谁能把 Coding/Agent 变成组织主线，谁就可能抓住第二幕；谁被上一阶段成功锁住，谁就可能掉队。

